\chapter{Benchmark 设计：输入、热路径和报告}

\section{本章要解决的问题}

上一章讲的是“为什么性能测量容易错”。本章讲“怎样把测量原则做成可复现的 benchmark 工程”。

这一步很关键。你以后做 HPC kernel、AI 算子、编译器优化、内存布局优化时，不可能每次都手写一段临时代码计时。你需要的是一套可以扩展、可以审计、可以保存原始数据、可以扫描参数、可以比较版本的实验系统。

本章目标是让你能设计一个小型 benchmark suite，而不是只会写一个 \code{main()} 计时器。这个 suite 至少要处理：

\begin{itemize}
  \item 多个 kernel：例如 \code{sum_u64}、\code{copy_u64}、\code{fill_u64}、\code{dot_f32}。
  \item 多个输入规模：从 L1 cache 级别到内存带宽级别。
  \item 多种数据分布：随机、全零、递增、分支友好、分支不友好。
  \item 多个重复次数：warmup、measured repetitions、batch。
  \item 正确性校验：每个 kernel 都有 reference 或可验证不变量。
  \item 原始数据和汇总数据：CSV 输出、环境记录、命令行记录。
  \item 二进制审计：可以把源码、反汇编和数字联系起来。
\end{itemize}

从工程角度看，benchmark suite 和普通测试框架很像：它要有清晰的数据流、明确的边界、稳定的输出格式和可维护的扩展点。差别在于，benchmark 还必须控制优化器、cache、频率、输入分布和统计口径。

\begin{keyidea}
一个好的 benchmark 不是一段计时代码，而是一套实验协议。它规定被测对象、输入空间、计时边界、重复策略、输出格式、正确性检查、环境记录和审计方法。
\end{keyidea}

\section{Benchmark 设计的主线}

把 benchmark 设计成下面的流水线：

\begin{lstlisting}
configuration:
    parse command line
    choose sizes, repetitions, batch, seed, output path

environment record:
    compiler flags, command line, CPU info, OS info, git commit if available

workload construction:
    allocate buffers
    initialize inputs
    compute reference or invariant

warmup:
    run kernel without recording samples

measurement:
    repeat:
        run timed region
        consume result
        validate if cheap enough
        append raw sample

summary:
    compute min, median, p95, max, derived metrics
    write CSV

audit:
    keep binary, objdump excerpt, perf command, raw data
\end{lstlisting}

这条流水线背后的思想是：每个阶段只做一类事。输入生成不放进计时区域；统计不混进 kernel；报告不改变运行行为；校验不依赖性能版本本身。

\begin{mentalmodel}
设计 benchmark 时，把它看成“实验编译器”：配置是源程序，workload 是输入实例，kernel 是被执行的目标代码，CSV 是输出产物，报告是解释器。任何隐式状态都会降低可复现性。
\end{mentalmodel}

\section{先写清楚问题陈述}

动手写代码前，先写一段 problem statement。它不需要长，但必须具体。

差的描述：

\begin{lstlisting}
measure sum performance
\end{lstlisting}

好的描述：

\begin{lstlisting}
Goal:
    Measure warm-cache and streaming performance of unsigned 64-bit sum.

Kernel:
    std::uint64_t sum_u64(std::span<const std::uint64_t> values)

Input sizes:
    1 KiB, 16 KiB, 256 KiB, 4 MiB, 64 MiB worth of uint64_t elements

Data:
    deterministic pseudo-random uint64_t, generated once per size

Measurement:
    5 warmup runs, 31 measured repetitions
    timed region contains only the kernel call

Output:
    raw CSV and summary CSV
    fields include n, bytes, sample_ns, median_ns, bandwidth_GBps
\end{lstlisting}

为什么要这么写？因为性能结论总是有适用范围。一个只在 \code{n = 1024} 时更快的版本，不一定在 \code{n = 16777216} 时更快；一个 warm-cache 下更快的版本，不一定 cold-cache 下更快；一个随机数据下快的分支版本，不一定排序数据下快。

\section{输入规模：不要只测一个点}

输入规模决定瓶颈。对数组 kernel 来说，常见阶段如下：

\begin{center}
\begin{tabular}{p{0.20\linewidth}p{0.32\linewidth}p{0.32\linewidth}}
\toprule
工作集范围 & 常见瓶颈 & 观察重点 \\
\midrule
很小 & 调用开销、循环开销、前端开销。 & ns/op、cycles/op。 \\
L1/L2 内 & 指令吞吐、load/store 吞吐、向量化。 & 每周期元素数。 \\
LLC 内 & cache 带宽、预取、并行度。 & GB/s、cache miss。 \\
超过 LLC & 内存带宽、地址转换和 NUMA。 & GB/s、LLC miss、TLB。 \\
\bottomrule
\end{tabular}
\end{center}

推荐选择不是随便的 \code{1000, 10000, 100000}，而是能跨过 cache 层级的规模。例如对 \code{std::uint64_t}：

\begin{lstlisting}
bytes per element = 8

sizes by bytes:
    1 KiB       -> 128 elements
    16 KiB      -> 2048 elements
    256 KiB     -> 32768 elements
    4 MiB       -> 524288 elements
    64 MiB      -> 8388608 elements
\end{lstlisting}

在命令行或配置中可以写成：

\begin{lstlisting}
--sizes-bytes 1024,16384,262144,4194304,67108864
\end{lstlisting}

还要测边界规模。比如 SIMD kernel 常见尾部处理问题：

\begin{lstlisting}
vector width = 8 float elements

important n:
    0, 1, 7, 8, 9, 15, 16, 17
    1023, 1024, 1025
\end{lstlisting}

这些规模不一定用来报告最终吞吐，但必须用于 correctness 和尾部路径性能检查。

\begin{warningbox}
只测一个大输入很容易掩盖小规模延迟；只测一个小输入又会掩盖内存带宽瓶颈。算子优化通常需要同时报告小规模、中规模、大规模和尾部边界。
\end{warningbox}

\section{数据分布：输入不是无关变量}

数据分布会改变分支、向量化、数值路径、cache 压缩效果和异常值。

以 \code{count_positive} 为例：

\begin{lstlisting}[language=C++]
std::uint64_t count_positive(std::span<const std::int32_t> values)
{
    std::uint64_t count = 0;
    for (const std::int32_t value : values) {
        if (value > 0) {
            ++count;
        }
    }
    return count;
}
\end{lstlisting}

这些输入会导致不同表现：

\begin{center}
\begin{tabular}{p{0.24\linewidth}p{0.46\linewidth}p{0.16\linewidth}}
\toprule
分布 & 特征 & 预测难度 \\
\midrule
全正数 & 分支总是 taken。 & 低 \\
全负数 & 分支总是 not taken。 & 低 \\
正负交替 & 模式简单但频繁变化。 & 中 \\
随机正负 & 近似不可预测。 & 高 \\
长 run & 一段正、一段负。 & 低到中 \\
\bottomrule
\end{tabular}
\end{center}

如果报告只写 \code{count_positive} 的速度，不写数据分布，结论不完整。

\subsection{确定性输入生成}

benchmark 不应该依赖每次运行都不同的随机输入。使用固定 seed：

\begin{lstlisting}[language=C++]
#include <cstdint>
#include <vector>

namespace {

constexpr std::uint64_t INPUT_MULTIPLIER = 6364136223846793005ULL;
constexpr std::uint64_t INPUT_INCREMENT = 1442695040888963407ULL;
constexpr std::uint64_t INPUT_DEFAULT_SEED = 1ULL;

std::vector<std::uint64_t> make_u64_input(std::size_t size, std::uint64_t seed)
{
    std::vector<std::uint64_t> values(size);
    std::uint64_t state = seed;
    for (std::uint64_t& value : values) {
        state = state * INPUT_MULTIPLIER + INPUT_INCREMENT;
        value = state;
    }
    return values;
}

}  // namespace
\end{lstlisting}

这不是密码学随机数，也不是统计质量最好的生成器。它的作用是：

\begin{itemize}
  \item 生成运行时数据，避免常量折叠。
  \item 固定 seed，保证可复现。
  \item 成本低，但仍然放在 setup，不放在 timed region。
\end{itemize}

\subsection{浮点数据要避开隐藏语义}

浮点 benchmark 要明确是否包含：

\begin{itemize}
  \item NaN。
  \item Inf。
  \item signed zero。
  \item denormal/subnormal。
  \item 极大或极小值。
\end{itemize}

例如 dot product 入门实验可以先使用 \code{[-1, 1]} 的普通有限值，避免把异常数值路径混进第一个实验。后续研究 denormal、fast-math、flush-to-zero 时，再单独设计对应 workload。

\section{对齐、步长和别名}

数组 kernel 的输入不只是长度，还有地址形状。

\begin{lstlisting}
contiguous:
    a[i]

strided:
    a[i * stride]

offset-aligned:
    pointer address is 64-byte aligned

misaligned:
    pointer address = aligned_base + 4
\end{lstlisting}

这些差异会影响：

\begin{itemize}
  \item load/store 是否跨 cache line。
  \item SIMD load 是否需要 unaligned 路径。
  \item 预取器能否识别访问模式。
  \item TLB 覆盖范围。
  \item 编译器能否证明无别名。
\end{itemize}

如果 kernel 接收两个指针：

\begin{lstlisting}[language=C++]
void copy_u64(std::span<std::uint64_t> dst,
              std::span<const std::uint64_t> src);
\end{lstlisting}

要明确：

\begin{itemize}
  \item \code{dst} 和 \code{src} 是否允许重叠。
  \item 若不允许重叠，C++ 实现是否能用 \code{restrict}-like 信息表达。
  \item 输入是否按 64 字节对齐。
  \item 是否测试非对齐偏移。
\end{itemize}

\begin{deepdive}
C++ 标准没有标准化的 \code{restrict} 关键字，但 GCC/Clang 有 \code{__restrict__} 扩展。后续讲 alias analysis 和 vectorization 时会看到：编译器是否知道两个指针不别名，可能直接决定是否能向量化。
\end{deepdive}

\section{热路径边界：把 setup 赶出去}

一个可维护 benchmark 应有显式阶段：

\begin{lstlisting}
make_workload()
compute_reference()
warmup()
measure()
validate()
write_result()
\end{lstlisting}

不要在 timed region 里做这些事：

\begin{itemize}
  \item 分配内存。
  \item 生成随机数。
  \item 打印日志。
  \item 写 CSV。
  \item 构造 \code{std::string}。
  \item 查询 CPU 信息。
  \item 解析命令行。
  \item 做复杂校验。
\end{itemize}

timed region 应尽量像这样：

\begin{lstlisting}[language=C++]
const auto begin = Clock::now();
const std::uint64_t result = sum_u64(workload.values);
const auto end = Clock::now();
do_not_optimize(result);
\end{lstlisting}

如果 kernel 本身就是“分配一个对象”或“写一个文件”，当然可以把分配或 I/O 放进 timed region。但那时 benchmark 名称要诚实：

\begin{lstlisting}
bench_vector_push_back_with_growth
bench_write_4KiB_sync
\end{lstlisting}

而不是假装它是纯计算 kernel。

\section{Batch：短 kernel 的必要技巧}

如果一次 kernel 调用只需几十纳秒，计时器开销和噪声会占很大比例。解决方法是 batch：

\begin{lstlisting}
one sample:
    begin timer
    repeat kernel B times
    end timer

sample_ns_per_batch = end - begin
ns_per_call = sample_ns_per_batch / B
\end{lstlisting}

但 batch 不是越大越好。它会改变 cache 状态、分支预测状态和热路径行为。比如同一个小数组重复求和很多次，会变成强 warm-cache 测试；如果你想测 cold-cache，batch 可能破坏实验。

一个简单策略：

\begin{itemize}
  \item 对短 kernel 使用 batch，目标是让每个 sample 至少达到几十微秒。
  \item 报告 batch 大小。
  \item 不同输入规模可以使用不同 batch，但要在 CSV 中记录。
  \item 不能把 batch 后的数字和非 batch 数字混在一起比较。
\end{itemize}

\begin{lstlisting}[language=C++]
std::uint64_t run_sum_batch(std::span<const std::uint64_t> values,
                            int batch_count)
{
    std::uint64_t combined = 0;
    for (int batch = 0; batch < batch_count; ++batch) {
        combined += sum_u64(values);
    }
    return combined;
}
\end{lstlisting}

\code{combined} 的作用是把每次结果连接起来，避免优化器删除重复调用。实际 harness 还要用 \code{do_not_optimize(combined)}。

\begin{warningbox}
batch 会改变 workload。对 cache 敏感、状态ful、随机输入、分支预测敏感的 kernel，batch 可能让测量更稳定，也可能让测量偏离真实场景。报告里必须写明。
\end{warningbox}

\section{CSV：先保存原始数据}

不要只打印最终汇总。原始样本是后续统计和复查的基础。

推荐至少输出两个文件：

\begin{lstlisting}
raw_samples.csv
summary.csv
\end{lstlisting}

raw samples：

\begin{lstlisting}
kernel,n,bytes,distribution,batch,rep,warmup,elapsed_ns,result_checksum
sum_u64,1048576,8388608,lcg,1,0,5,312000,12345
sum_u64,1048576,8388608,lcg,1,1,5,311400,12345
\end{lstlisting}

summary：

\begin{lstlisting}
kernel,n,bytes,distribution,batch,repetitions,min_ns,median_ns,p95_ns,max_ns,GBps
sum_u64,1048576,8388608,lcg,1,31,309900,312000,320100,330000,26.88
\end{lstlisting}

字段设计原则：

\begin{itemize}
  \item 每一行都能独立解释。
  \item 不要把单位藏在列名之外。
  \item 所有影响结果的参数都要进 CSV。
  \item 原始样本不要覆盖，只追加或写到带时间戳/配置名的文件。
  \item summary 可以由 raw samples 再生成。
\end{itemize}

\section{环境记录：报告不是记忆}

benchmark 输出目录建议包含：

\begin{lstlisting}
results/2026-06-13_sum_suite/
    raw_samples.csv
    summary.csv
    environment.txt
    command.txt
    objdump.txt
    perf_stat.txt
    report.md
\end{lstlisting}

\code{environment.txt} 至少包含：

\begin{lstlisting}[language=bash]
date -Iseconds
uname -a
lscpu
clang++ --version
g++ --version
cat /sys/devices/system/cpu/cpu0/cpufreq/scaling_governor
cat /proc/sys/kernel/perf_event_paranoid
\end{lstlisting}

这些命令不一定每台机器都有同样输出。没关系，记录能记录的，并把失败也记录下来：

\begin{lstlisting}
scaling_governor:
    unavailable: /sys/devices/system/cpu/cpu0/cpufreq does not exist
\end{lstlisting}

这比假装环境不存在要可靠。

\section{一个小型 benchmark suite 的数据模型}

先设计数据结构。不要一开始就把所有逻辑塞进 \code{main()}。

\begin{lstlisting}[language=C++]
#include <cstdint>
#include <string>
#include <vector>

struct BenchmarkOptions {
    std::vector<std::size_t> sizes;
    int warmup_repetitions = 0;
    int measured_repetitions = 0;
    int batch_count = 0;
    std::uint64_t seed = 0;
    std::string output_prefix;
};

struct RawSample {
    std::string kernel_name;
    std::size_t elements = 0;
    std::size_t bytes = 0;
    int batch_count = 0;
    int repetition = 0;
    std::uint64_t elapsed_ns = 0;
    std::uint64_t checksum = 0;
};

struct SummaryRow {
    std::string kernel_name;
    std::size_t elements = 0;
    std::size_t bytes = 0;
    int batch_count = 0;
    int repetitions = 0;
    std::uint64_t min_ns = 0;
    std::uint64_t median_ns = 0;
    std::uint64_t p95_ns = 0;
    std::uint64_t max_ns = 0;
    double gb_per_second = 0.0;
};
\end{lstlisting}

这些结构分别对应：

\begin{itemize}
  \item \code{BenchmarkOptions}：实验配置。
  \item \code{RawSample}：每次测量的原始样本。
  \item \code{SummaryRow}：由原始样本计算出的汇总。
\end{itemize}

这里先用 public data struct，因为它们只是数据载体。如果后续需要不变量检查、解析和格式化，可以再封装成类。

\section{命令行参数：可复现需要显式配置}

硬编码常量适合第一章入门，但 benchmark suite 应该能从命令行配置。

示例：

\begin{lstlisting}[language=bash]
./bench_suite \
    --sizes 1024,16384,262144,1048576 \
    --warmup 5 \
    --repetitions 31 \
    --batch 1 \
    --seed 1 \
    --output results/sum_suite
\end{lstlisting}

解析器可以先写得简单，但要有清晰错误：

\begin{lstlisting}[language=C++]
#include <charconv>
#include <cstdint>
#include <cstdlib>
#include <string_view>

namespace {

bool parse_u64(std::string_view text, std::uint64_t* value_out)
{
    const char* first = text.data();
    const char* last = text.data() + text.size();
    std::uint64_t value = 0;
    const auto result = std::from_chars(first, last, value);
    if (result.ec != std::errc{} || result.ptr != last) {
        return false;
    }
    *value_out = value;
    return true;
}

}  // namespace
\end{lstlisting}

不要在解析失败后悄悄用默认值。性能报告里的配置必须可信；输入错了就应当失败并打印错误。

\begin{deepdive}
成熟项目可以使用 CLI11、cxxopts、absl flags 等库。本书早期用手写解析器，是为了让你看清楚参数如何进入实验系统。工程项目中应根据规模选择库，避免自制复杂解析器。
\end{deepdive}

\section{Workload factory：生成输入但不污染计时}

对不同 kernel，可以定义不同 workload。

\begin{lstlisting}[language=C++]
#include <cstdint>
#include <span>
#include <vector>

struct SumWorkload {
    std::vector<std::uint64_t> values;
    std::uint64_t expected = 0;
};

struct CopyWorkload {
    std::vector<std::uint64_t> source;
    std::vector<std::uint64_t> destination;
};

struct DotWorkload {
    std::vector<float> a;
    std::vector<float> b;
    float expected = 0.0f;
};
\end{lstlisting}

factory 负责：

\begin{itemize}
  \item 分配 buffer。
  \item 初始化输入。
  \item 计算 reference。
  \item 保证数据在 timed region 前准备好。
\end{itemize}

\begin{lstlisting}[language=C++]
std::uint64_t sum_u64_ref(std::span<const std::uint64_t> values)
{
    std::uint64_t sum = 0;
    for (const std::uint64_t value : values) {
        sum += value;
    }
    return sum;
}

SumWorkload make_sum_workload(std::size_t elements, std::uint64_t seed)
{
    SumWorkload workload;
    workload.values = make_u64_input(elements, seed);
    workload.expected = sum_u64_ref(workload.values);
    return workload;
}
\end{lstlisting}

对 \code{copy_u64}，reference 可以不是另一个 copy 函数，而是校验不变量：

\begin{lstlisting}
after copy:
    destination[i] == source[i] for all i
\end{lstlisting}

对 \code{fill_u64}：

\begin{lstlisting}
after fill:
    destination[i] == fill_value for all i
\end{lstlisting}

对 \code{dot_f32}，reference 需要误差容忍：

\begin{lstlisting}
abs(actual - expected) <= epsilon * max(abs(expected), 1.0)
\end{lstlisting}

\section{Kernel API：统一但不过度抽象}

为了让 suite 能跑多个 kernel，需要统一接口。但不要为了统一而隐藏重要差异。

可以从简单函数开始：

\begin{lstlisting}[language=C++]
std::uint64_t sum_u64_kernel(std::span<const std::uint64_t> values)
{
    std::uint64_t sum = 0;
    for (const std::uint64_t value : values) {
        sum += value;
    }
    return sum;
}

void copy_u64_kernel(std::span<std::uint64_t> destination,
                     std::span<const std::uint64_t> source)
{
    for (std::size_t i = 0; i < source.size(); ++i) {
        destination[i] = source[i];
    }
}

void fill_u64_kernel(std::span<std::uint64_t> destination,
                     std::uint64_t value)
{
    for (std::uint64_t& element : destination) {
        element = value;
    }
}

float dot_f32_kernel(std::span<const float> a, std::span<const float> b)
{
    float sum = 0.0f;
    for (std::size_t i = 0; i < a.size(); ++i) {
        sum += a[i] * b[i];
    }
    return sum;
}
\end{lstlisting}

这里没有强行把所有 kernel 包成同一个函数指针类型，因为它们返回值和参数不同。更实际的做法是每类 kernel 有自己的 runner，然后输出相同的 \code{RawSample}。

\begin{warningbox}
过度抽象会污染 benchmark。虚函数调用、\code{std::function} type erasure、动态分派、分配和锁都可能进入热路径。抽象应放在 timed region 外，或者确认编译器能完全内联。
\end{warningbox}

\section{校验策略：不同 kernel 不同方法}

正确性校验不一定每次都完整扫描，但必须有明确策略。

\begin{center}
\begin{tabular}{p{0.20\linewidth}p{0.30\linewidth}p{0.34\linewidth}}
\toprule
kernel & 校验方式 & 注意 \\
\midrule
\code{sum_u64} & 比较返回值。 & 使用 unsigned 避免 signed overflow UB。 \\
\code{copy_u64} & 比较输出和输入。 & 校验会扫描内存，不放入 timed region。 \\
\code{fill_u64} & 检查元素等于常量。 & 测量前可能要重置 buffer。 \\
\code{dot_f32} & 误差容忍比较。 & 说明容忍度和 reference 顺序。 \\
\bottomrule
\end{tabular}
\end{center}

对于会修改输出 buffer 的 kernel，要注意每次测量前的初始状态。比如 \code{fill_u64}：

\begin{lstlisting}
for each repetition:
    reset destination outside timed region if required
    time fill_u64(destination, value)
    validate destination outside timed region
\end{lstlisting}

如果 reset 成本很大，不能混进 timed region；如果不 reset，又要确认 kernel 的性能不依赖原始内容。

\section{样本汇总：从 raw 到 summary}

summary 应由 raw samples 计算。最小实现：

\begin{lstlisting}[language=C++]
#include <algorithm>
#include <cstdint>
#include <vector>

std::uint64_t percentile_sorted(const std::vector<std::uint64_t>& sorted,
                                std::size_t numerator,
                                std::size_t denominator)
{
    if (sorted.empty()) {
        return 0;
    }
    const std::size_t last_index = sorted.size() - 1;
    const std::size_t index = (last_index * numerator + denominator - 1) / denominator;
    return sorted[index];
}

SummaryRow summarize_samples(std::string kernel_name,
                             std::size_t elements,
                             std::size_t bytes,
                             int batch_count,
                             std::vector<std::uint64_t> elapsed_ns)
{
    constexpr std::size_t P95_NUMERATOR = 95;
    constexpr std::size_t PERCENT_DENOMINATOR = 100;
    constexpr double BYTES_PER_GIGABYTE = 1.0e9;
    constexpr double NANOSECONDS_PER_SECOND = 1.0e9;

    std::sort(elapsed_ns.begin(), elapsed_ns.end());

    SummaryRow row;
    row.kernel_name = std::move(kernel_name);
    row.elements = elements;
    row.bytes = bytes;
    row.batch_count = batch_count;
    row.repetitions = static_cast<int>(elapsed_ns.size());
    row.min_ns = elapsed_ns.front();
    row.median_ns = elapsed_ns[elapsed_ns.size() / 2];
    row.p95_ns = percentile_sorted(elapsed_ns, P95_NUMERATOR, PERCENT_DENOMINATOR);
    row.max_ns = elapsed_ns.back();

    const double seconds = static_cast<double>(row.median_ns) / NANOSECONDS_PER_SECOND;
    row.gb_per_second = static_cast<double>(bytes) / seconds / BYTES_PER_GIGABYTE;
    return row;
}
\end{lstlisting}

这个实现还有局限：

\begin{itemize}
  \item 空样本只返回 0，生产代码应报告错误。
  \item p95 使用简单 nearest-rank 近似。
  \item GB/s 默认每个 sample 只处理一次 \code{bytes}，如果 batch 大于 1，公式要乘以 batch。
\end{itemize}

所以更准确：

\begin{lstlisting}
effective_bytes = bytes * batch_count
GB/s = effective_bytes / seconds / 1e9
\end{lstlisting}

\section{派生指标：不要只给时间}

不同 kernel 应输出不同派生指标。

\subsection{\code{sum_u64}}

显式读：

\begin{lstlisting}
bytes_read = n * sizeof(uint64_t)
GB/s = bytes_read * batch / seconds / 1e9
ns/element = elapsed_ns / (n * batch)
\end{lstlisting}

\subsection{\code{copy_u64}}

copy 既读又写：

\begin{lstlisting}
bytes_read = n * sizeof(uint64_t)
bytes_written = n * sizeof(uint64_t)
logical_bytes = bytes_read + bytes_written
GB/s = logical_bytes * batch / seconds / 1e9
\end{lstlisting}

硬件实际流量可能更高，因为写 miss 可能触发 write allocate。后续讲内存系统时会展开。

\subsection{\code{fill_u64}}

fill 主要写：

\begin{lstlisting}
bytes_written = n * sizeof(uint64_t)
GB/s = bytes_written * batch / seconds / 1e9
\end{lstlisting}

但普通 store 可能也触发 write allocate。非时序 store 是另一个话题。

\subsection{\code{dot_f32}}

每个元素通常算 1 次乘法和 1 次加法：

\begin{lstlisting}
FLOPs = 2 * n
GFLOP/s = FLOPs * batch / seconds / 1e9
bytes_read = 2 * n * sizeof(float)
\end{lstlisting}

如果使用 FMA，仍通常按 2 FLOPs 计。

\section{输出 CSV 的代码结构}

不要在测量循环里频繁 flush。先收集样本，最后写文件。

\begin{lstlisting}[language=C++]
#include <cstdio>
#include <string>
#include <vector>

void write_raw_samples_csv(const std::string& path,
                           const std::vector<RawSample>& samples)
{
    FILE* file = std::fopen(path.c_str(), "w");
    if (file == nullptr) {
        std::perror("fopen raw samples");
        std::abort();
    }

    std::fprintf(file,
                 "kernel,n,bytes,batch,rep,elapsed_ns,checksum\n");

    for (const RawSample& sample : samples) {
        std::fprintf(file,
                     "%s,%zu,%zu,%d,%d,%llu,%llu\n",
                     sample.kernel_name.c_str(),
                     sample.elements,
                     sample.bytes,
                     sample.batch_count,
                     sample.repetition,
                     static_cast<unsigned long long>(sample.elapsed_ns),
                     static_cast<unsigned long long>(sample.checksum));
    }

    std::fclose(file);
}
\end{lstlisting}

这段代码使用 C stdio 是为了简单直接。生产代码可以用 RAII 封装文件句柄，或者使用 \code{std::ofstream}。关键不是 API 选择，而是：文件写入不在 timed region 里。

\section{Runner：把一次实验组织起来}

下面是一个简化的 sum runner。它展示结构，不追求覆盖所有错误处理。

\begin{lstlisting}[language=C++]
#include <chrono>
#include <cstdint>
#include <span>
#include <vector>

namespace {

using Clock = std::chrono::steady_clock;

template <class T>
inline void do_not_optimize(const T& value)
{
    asm volatile("" : : "g"(value) : "memory");
}

std::uint64_t elapsed_ns_since(Clock::time_point begin, Clock::time_point end)
{
    return static_cast<std::uint64_t>(
        std::chrono::duration_cast<std::chrono::nanoseconds>(end - begin).count());
}

std::uint64_t run_sum_batch(std::span<const std::uint64_t> values,
                            int batch_count)
{
    std::uint64_t combined = 0;
    for (int batch = 0; batch < batch_count; ++batch) {
        combined += sum_u64_kernel(values);
    }
    return combined;
}

void warmup_sum(const SumWorkload& workload,
                int warmup_repetitions,
                int batch_count)
{
    for (int iter = 0; iter < warmup_repetitions; ++iter) {
        const std::uint64_t result = run_sum_batch(workload.values, batch_count);
        do_not_optimize(result);
    }
}

RawSample measure_sum_once(const SumWorkload& workload,
                           std::size_t elements,
                           int batch_count,
                           int repetition)
{
    const auto begin = Clock::now();
    const std::uint64_t result = run_sum_batch(workload.values, batch_count);
    const auto end = Clock::now();

    do_not_optimize(result);

    RawSample sample;
    sample.kernel_name = "sum_u64";
    sample.elements = elements;
    sample.bytes = elements * sizeof(std::uint64_t);
    sample.batch_count = batch_count;
    sample.repetition = repetition;
    sample.elapsed_ns = elapsed_ns_since(begin, end);
    sample.checksum = result;
    return sample;
}

}  // namespace
\end{lstlisting}

注意：

\begin{itemize}
  \item warmup 和 measurement 分开。
  \item timed region 只包 batch kernel。
  \item \code{RawSample} 记录 batch 和 checksum。
  \item 校验可以在外层比较 checksum 和 expected batch result。
\end{itemize}

expected batch result：

\begin{lstlisting}
expected_batch = workload.expected * batch_count
\end{lstlisting}

因为使用 \code{std::uint64_t}，乘法和加法都按模 \code{2^64} 回绕，语义明确。

\section{参数扫描：矩阵比单点更重要}

真实 benchmark 通常是多维矩阵：

\begin{lstlisting}
kernel:
    sum_u64, copy_u64, fill_u64, dot_f32

size:
    1 KiB, 16 KiB, 256 KiB, 4 MiB, 64 MiB

distribution:
    zero, lcg, increasing

cache mode:
    warm, cold-ish

batch:
    1, auto
\end{lstlisting}

不要一开始就跑满所有组合。组合爆炸会浪费时间。训练建议：

\begin{enumerate}
  \item 先固定 distribution 和 cache mode，扫描 size。
  \item 发现异常后，增加 distribution 维度。
  \item 对短 kernel 引入 batch。
  \item 对内存 kernel 引入 cold-ish 模式。
  \item 最后再做完整矩阵。
\end{enumerate}

CSV 里必须记录每个维度，否则后续无法比较。

\section{自动 batch 的思路}

自动 batch 的目标是让单个 sample 足够长。例如希望每个 sample 至少 100 微秒：

\begin{lstlisting}
target_sample_ns = 100000
batch = 1
while measured_ns(batch) < target_sample_ns:
    batch *= 2
\end{lstlisting}

但这个校准本身也有成本和噪声。它应该在正式测量前完成，并记录最终 batch。

\begin{lstlisting}[language=C++]
int choose_batch_count(std::span<const std::uint64_t> values,
                       std::uint64_t target_ns)
{
    constexpr int BENCHMARK_MAX_BATCH = 1 << 20;

    int batch_count = 1;
    while (batch_count < BENCHMARK_MAX_BATCH) {
        const auto begin = Clock::now();
        const std::uint64_t result = run_sum_batch(values, batch_count);
        const auto end = Clock::now();
        do_not_optimize(result);

        if (elapsed_ns_since(begin, end) >= target_ns) {
            break;
        }
        batch_count *= 2;
    }
    return batch_count;
}
\end{lstlisting}

风险：

\begin{itemize}
  \item 校准时的数据会变热。
  \item batch 变大后测的是重复运行场景。
  \item 不同机器上 batch 不同，比较时要记录。
\end{itemize}

\section{避免 harness 干扰热路径}

benchmark framework 本身也会影响结果。常见干扰：

\begin{itemize}
  \item 使用 \code{std::function} 包装 kernel，导致间接调用。
  \item 在 timed region 里捕获 lambda，编译器无法内联。
  \item 每次 sample 创建临时 vector。
  \item 每次 sample 写日志或 CSV。
  \item 校验代码意外进入 timed region。
  \item runner 过度泛型，导致编译器生成复杂控制流。
\end{itemize}

对热路径最稳的办法是：

\begin{lstlisting}
outside timed region:
    choose kernel
    construct workload
    prepare function-specific runner

inside timed region:
    direct call or fully inlined call
\end{lstlisting}

如果使用函数指针：

\begin{lstlisting}[language=C++]
using SumKernel = std::uint64_t (*)(std::span<const std::uint64_t>);
\end{lstlisting}

函数指针调用可能不能内联，但可以模拟真实动态 dispatch。如果你的生产系统就是通过函数指针做 ISA dispatch，这可能是合理的。如果目标是测 kernel 本体上限，应尽量测 direct call 或单独测 dispatch overhead。

\begin{deepdive}
AI 算子库常有 dispatcher：根据 dtype、shape、ISA、线程数选择 kernel。优化时要分清两个 benchmark：一个测 end-to-end operator dispatch，一个测已经选定的 microkernel。二者都有价值，但不能混为一谈。
\end{deepdive}

\section{四个基础 kernel 的设计}

\subsection{\code{sum_u64}}

目的：

\begin{itemize}
  \item 训练只读 streaming kernel。
  \item 观察 cache 层级和内存带宽。
  \item 学习 DCE 防护和 checksum。
\end{itemize}

语义：

\begin{lstlisting}
sum = 0
for x in values:
    sum += x
return sum
\end{lstlisting}

指标：

\begin{lstlisting}
GB/s = n * sizeof(uint64_t) * batch / seconds / 1e9
ns/element = elapsed_ns / (n * batch)
\end{lstlisting}

\subsection{\code{copy_u64}}

目的：

\begin{itemize}
  \item 训练读加写 memory kernel。
  \item 观察 write allocate、带宽上限和 \code{memcpy} 对比。
  \item 学习输出 buffer 校验。
\end{itemize}

语义：

\begin{lstlisting}
for i in [0, n):
    dst[i] = src[i]
\end{lstlisting}

指标：

\begin{lstlisting}
logical bytes = 2 * n * sizeof(uint64_t)
\end{lstlisting}

\subsection{\code{fill_u64}}

目的：

\begin{itemize}
  \item 训练写密集 kernel。
  \item 观察 store 带宽、write allocate 和初始化成本。
  \item 后续可对比普通 store 与 non-temporal store。
\end{itemize}

语义：

\begin{lstlisting}
for i in [0, n):
    dst[i] = value
\end{lstlisting}

\subsection{\code{dot_f32}}

目的：

\begin{itemize}
  \item 训练浮点计算和内存读取混合 kernel。
  \item 为后续 SIMD、FMA、reduction 打基础。
  \item 学习 GFLOP/s 和数值误差报告。
\end{itemize}

语义：

\begin{lstlisting}
sum = 0
for i in [0, n):
    sum += a[i] * b[i]
return sum
\end{lstlisting}

指标：

\begin{lstlisting}
FLOPs = 2 * n
bytes_read = 2 * n * sizeof(float)
\end{lstlisting}

\section{报告中的图表和文字}

就算没有画图，也要让表格能表达趋势。推荐 summary 表：

\begin{lstlisting}
kernel,n,bytes,median_ns,GBps,ns_per_element,notes
sum_u64,128,1024,80,12.8,0.625,L1-sized
sum_u64,32768,262144,9000,29.1,0.275,L2-sized
sum_u64,8388608,67108864,3500000,19.2,0.417,memory-sized
\end{lstlisting}

报告文字不要只写“随着 n 增大性能下降”。要写因果假设：

\begin{lstlisting}
Observation:
    Throughput rises from 1 KiB to 256 KiB, then drops at 64 MiB.

Hypothesis:
    Small sizes are dominated by loop and timer overhead.
    Middle sizes benefit from cache residency.
    Large sizes are limited by memory bandwidth and TLB pressure.

Evidence:
    objdump shows scalar loop, no SIMD.
    perf stat shows cache misses increase at large n.

Limit:
    CPU frequency was not pinned; result needs repeat under performance governor.
\end{lstlisting}

\section{实验 16.1：设计输入规模扫描}

\begin{labbox}
目标：把单点 benchmark 改成输入规模扫描。

基于第 15 章的 \code{sum_u64} benchmark，支持下面命令：

\begin{lstlisting}[language=bash]
./bench_sum_sizes --sizes 128,2048,32768,524288,8388608 \
                  --warmup 5 \
                  --repetitions 31 \
                  --seed 1 \
                  --output results/sum_sizes
\end{lstlisting}

任务：

\begin{enumerate}
  \item 实现 \code{--sizes} 解析，逗号分隔。
  \item 每个 size 单独构造 workload。
  \item 每个 size 单独 warmup 和测量。
  \item 输出 \code{raw_samples.csv} 和 \code{summary.csv}。
  \item summary 中加入 \code{GBps} 和 \code{ns_per_element}。
\end{enumerate}

交付物：

\begin{enumerate}
  \item 源码。
  \item 命令行。
  \item raw CSV 前 10 行。
  \item summary CSV。
  \item 一段趋势解释。
\end{enumerate}
\end{labbox}

\section{实验 16.2：加入 \code{copy_u64} 和 \code{fill_u64}}

\begin{labbox}
目标：把单 kernel benchmark 扩展成多 kernel suite。

新增：

\begin{lstlisting}[language=C++]
void copy_u64_kernel(std::span<std::uint64_t> dst,
                     std::span<const std::uint64_t> src);

void fill_u64_kernel(std::span<std::uint64_t> dst,
                     std::uint64_t value);
\end{lstlisting}

任务：

\begin{enumerate}
  \item 为 \code{copy_u64} 设计 workload。
  \item 为 \code{fill_u64} 设计 workload。
  \item 分别定义 correctness check。
  \item CSV 中增加 \code{kernel} 列。
  \item 对每个 kernel 正确计算 logical bytes。
  \item 对比 \code{copy_u64} 和 \code{std::memcpy}，但要分开报告。
\end{enumerate}

交付物：

\begin{enumerate}
  \item 扩展后的 suite 源码。
  \item 三个 kernel 的 summary。
  \item 每个 kernel 的指标定义。
  \item \code{copy_u64} 与 \code{std::memcpy} 的对比报告。
\end{enumerate}
\end{labbox}

\section{实验 16.3：数据分布对分支的影响}

\begin{labbox}
目标：观察输入分布如何改变分支 kernel 性能。

实现：

\begin{lstlisting}[language=C++]
std::uint64_t count_positive(std::span<const std::int32_t> values);
\end{lstlisting}

生成 4 种输入：

\begin{itemize}
  \item 全正数。
  \item 全负数。
  \item 正负交替。
  \item 固定 seed 的随机正负。
\end{itemize}

任务：

\begin{enumerate}
  \item 每种分布使用同样 size、warmup、repetitions。
  \item CSV 中记录 \code{distribution}。
  \item 用 \code{perf stat} 记录 \code{branches} 和 \code{branch-misses}。
  \item 反汇编核心循环，说明是否使用分支、\instruction{cmov} 或 \instruction{setcc}。
  \item 解释性能差异。
\end{enumerate}

交付物：

\begin{enumerate}
  \item 四种分布的结果表。
  \item \code{perf stat} 输出。
  \item 核心循环反汇编。
  \item 800 字报告：输入分布、预测器和性能之间的关系。
\end{enumerate}
\end{labbox}

\section{实验 16.4：batch 校准}

\begin{labbox}
目标：为短 kernel 设计自动 batch，并分析它的副作用。

任务：

\begin{enumerate}
  \item 实现 \code{--batch auto}。
  \item 设定目标单样本时间，例如 100 微秒。
  \item 对每个 size 选择 batch。
  \item CSV 记录最终 batch。
  \item 比较固定 batch 1 和 auto batch 的结果稳定性。
\end{enumerate}

交付物：

\begin{enumerate}
  \item batch 校准代码。
  \item 每个 size 的 batch 表。
  \item auto batch 前后的 min、median、max。
  \item 说明 batch 是否改变了 cache 状态。
\end{enumerate}
\end{labbox}

\section{实验 16.5：完整 benchmark 报告目录}

\begin{labbox}
目标：生成一个可以提交给别人复现的 benchmark 结果目录。

目录：

\begin{lstlisting}
results/ch16_suite/
    command.txt
    environment.txt
    raw_samples.csv
    summary.csv
    objdump.txt
    perf_stat.txt
    report.md
\end{lstlisting}

任务：

\begin{enumerate}
  \item 运行 suite，并保存完整命令到 \code{command.txt}。
  \item 保存 \code{lscpu}、\code{uname -a}、编译器版本。
  \item 保存 \code{objdump -drwC -Mintel} 输出。
  \item 对至少一个 kernel 运行 \code{perf stat}。
  \item 写 \code{report.md}，包含目标、方法、结果、解释和限制。
\end{enumerate}

交付物：

\begin{enumerate}
  \item 完整结果目录。
  \item 一份不超过 1500 字但信息完整的报告。
  \item 至少 3 条你认为需要进一步实验验证的假设。
\end{enumerate}
\end{labbox}

\section{本章作业}

\begin{exercisebox}
\subsection*{作业 1：设计一个 benchmark suite 架构}

画出你的 benchmark suite 模块图，至少包含：

\begin{itemize}
  \item options parsing。
  \item workload factory。
  \item kernel runner。
  \item correctness checker。
  \item sample collector。
  \item summary generator。
  \item CSV writer。
  \item environment recorder。
\end{itemize}

要求说明每个模块输入、输出和不应承担的责任。

\subsection*{作业 2：输入规模和 cache 假设}

选择你的机器，根据 \code{lscpu} 或系统信息估计 L1/L2/L3 cache 大小。为 \code{sum_u64} 设计 8 个输入规模：

\begin{enumerate}
  \item 至少 2 个明显小于 L1。
  \item 至少 2 个接近 L2 或 L3。
  \item 至少 2 个明显超过 LLC。
  \item 至少 2 个尾部边界规模，例如非 8、16、32 的倍数。
\end{enumerate}

写出每个规模的元素数、字节数、预期瓶颈和验证方法。

\subsection*{作业 3：CSV schema 评审}

设计 raw 和 summary 两个 CSV schema。要求：

\begin{itemize}
  \item 每列都有单位或明确语义。
  \item 每个影响结果的参数都能记录。
  \item 能支持多 kernel、多 size、多 distribution、多 batch。
  \item 能从 raw 重新生成 summary。
\end{itemize}

提交 schema 和 5 行示例数据，并解释为什么这些列足够复现实验。

\subsection*{作业 4：多 kernel suite 实现}

实现一个 C++20 benchmark suite，至少包含：

\begin{itemize}
  \item \code{sum_u64}。
  \item \code{copy_u64}。
  \item \code{fill_u64}。
  \item \code{dot_f32}。
\end{itemize}

要求：

\begin{enumerate}
  \item 命令行可配置 size、warmup、repetitions、batch、seed。
  \item 每个 kernel 都有 correctness check。
  \item 输出 raw 和 summary CSV。
  \item 不把 allocation、random generation、CSV writing 放进 timed region。
  \item 提供至少一个 \code{objdump} 核心循环分析。
\end{enumerate}

\subsection*{作业 5：Benchmark 报告}

用你的 suite 做一次完整实验，提交报告：

\begin{itemize}
  \item 环境。
  \item 编译命令。
  \item workload matrix。
  \item raw data 路径。
  \item summary 表。
  \item 至少 2 个 kernel 的趋势解释。
  \item 至少 1 个 PMU 证据。
  \item 至少 1 个反汇编证据。
  \item 实验限制。
\end{itemize}

\subsection*{评分标准}

本章作业按下面标准评价：

\begin{enumerate}
  \item benchmark 架构是否职责清晰。
  \item 输入规模是否覆盖不同性能区域。
  \item 数据分布、batch、cache 状态是否记录。
  \item timed region 是否干净。
  \item CSV 是否能支持复现和后续统计。
  \item correctness check 是否独立可信。
  \item 报告是否用数据、汇编和硬件事件支撑结论。
\end{enumerate}
\end{exercisebox}

\section{本章小结}

本章把可信测量原则工程化成 benchmark suite 设计：

\begin{lstlisting}
benchmark suite:
    options
    workload factory
    kernel runner
    correctness checker
    sample collector
    summary generator
    CSV writer
    environment recorder

input matrix:
    size
    distribution
    alignment
    cache mode
    batch

outputs:
    raw_samples.csv
    summary.csv
    environment.txt
    command.txt
    objdump.txt
    perf_stat.txt
    report.md
\end{lstlisting}

到这里，你应该能把一个临时代码片段升级成可复现的实验系统。下一章会继续补齐统计与噪声分析：如何理解样本分布、异常值、均值和中位数，如何比较两个版本是否真的不同，以及如何把实验复现做得更严谨。
